\documentclass[10pt,a4paper]{amsart}
\usepackage[margin=19mm]{geometry}
\usepackage{amsmath,amssymb,mathtools}
\usepackage[scr=rsfs,cal=euler]{mathalfa}
\usepackage{xfrac}
\usepackage[unicode,hidelinks,bookmarks=false]{hyperref}
\newcommand{\ie}{i.e.\ }
\newcommand{\cf}{cf.\ }
\newcommand{\resp}{respectively\ }
\newcommand{\Subsection}{\subsection}
\providecommand{\nobreakdash}{\nobreak\hbox{-}}
\setlength{\emergencystretch}{2em}
\multlinegap0pt
\usepackage{amssymb}
\usepackage{mathtools}
\usepackage[shortlabels]{enumitem}
\usepackage{tikz}
\usepackage{bbm}
\usepackage{dsfont}
\newtheorem{proposition}{Proposition}
\newtheorem{lemma}{Lemma}
\renewcommand\ge{\geqslant}
\renewcommand\le{\leqslant}
\title{Colour diversity in spanning structures under Dirac-type conditions}
\author{Xinbu Cheng, Xinqi Huang, Hong Liu, Bin Wang, Zhifei Yan}
\date{Source: arXiv:2602.23801v1 (CC BY 4.0)}
\begin{document}
\maketitle

\noindent\emph{Source and licence.} Colour diversity in spanning structures under Dirac-type conditions. Version arXiv:2602.23801v1 (CC BY 4.0), distributed by arXiv under the Creative Commons Attribution 4.0 International licence, http://creativecommons.org/licenses/by/4.0/, as recorded in the arXiv licence metadata for this exact version and shown on the abstract page https://arxiv.org/abs/2602.23801v1. Full licence notice: \texttt{licenses/arxiv-2602.23801-CC-BY-4.0.txt}.\par\smallskip

\noindent\textit{Editorial scope.} Counted unit: the article's proposition prop:pathforest (source0.tex lines 325--327). The article's complete original proof of it is delivered with it (source0.tex lines 472--550, one \texttt{proof} environment of 79 lines, which the article places away from the statement). No mathematics is written, repaired or re-derived: the statement and the proof are the article's own lines, in the article's own notation, and no retained line is substituted or re-flowed.  Retained uncounted as the source-local context the proof needs: lines 361--368; lines 388--397.  Nothing was omitted from any retained span, so the package carries no omission record.  No retained line contains a citation, so the wrapper carries no bibliography and no bibliography entry.  No retained line contains a figure environment, an image-inclusion command or drawing code, so no image is delivered and the package declares no asset dependency.  Editorial formatting only, in addition: an AMS \texttt{amsart} preamble that restates the article's own declarations and macros used by the retained lines, namely the environments \texttt{proposition}, \texttt{lemma} on separate counters, together with the standard packages those lines need.  The retained environments print in this document as Proposition 1, Lemma 1 in order of appearance, whereas the article prints the same items with the numbering of its own full document.  The complete native source of the article as distributed in the arXiv e-print for this exact version is bundled unchanged as \texttt{sources/195422/source0.tex}.
\begin{equation}\label{eq:colourset}
C_i :=
\Big\{ f(x)\;:\;
x\in N_{C_{i-1}}\!\big(v_1^{(i)}\big)
      \setminus \{v_1^{(1)},\ldots,v_1^{(t)}\}
\Big\}
\;\cup\; C_{i-1},
\end{equation}

\begin{lemma}\label{lem:keyobs}
Let $F_0\in\mathcal F$ and $C_i$, $i=0,1,\ldots,t$ be defined as above. 
Suppose that the number of colours in $F_0$ is less than $(c-\delta)n$.
Then for every $1\le i\le t$ and every
\[
x\in N_{C_{i-1}}\!\big(v_1^{(i)}\big)\setminus \{v_1^{(1)},\ldots,v_1^{(t)}\} \quad\text{and}\quad
y\in V(F_0)\setminus\{x,v_1^{(1)},\ldots,v_1^{(t)}\},
\]
we have $f(x)\neq f(y)$.
\end{lemma}

\begin{proposition}\label{prop:pathforest}
For $1/2<c\le 1$, there exists $n_0\in\mathbb N$ such that for all $n\ge n_0$ the following holds. Let $\delta=\delta(n)$ and $t=t(n)$ satisfy $2c<t\delta<\delta^2 n/2$. If $G$ is a graph on $n$ vertices with $\delta(G)\ge cn$ and $\chi$ is a proper edge-colouring of $G$, then $G$ contains a spanning linear forest $F$ consisting of $t$ vertex-disjoint paths and containing at least $(c-\delta)n$ distinct colours.
\end{proposition}

\begin{proof}[Proof of Proposition~\ref{prop:pathforest}]
Let $F_0=\bigcup_{i=1}^t P_i\in\mathcal F$ be a spanning linear forest with exactly $t$
components that maximises the number of distinct colours, and suppose for
contradiction that $F_0$ uses fewer than $(c-\delta)n$ colours.
Let $C_0\subset C_1\subset\cdots\subset C_t$ be the colour sets defined in
\eqref{eq:colourset}, and let $C(F_0)=\chi(E(F_0))$ denote the set of colours appearing
in $F_0$.
Since $C_t\setminus C_0\subseteq C(F_0)$, it suffices to show that
\[
|C_t\setminus C_0|>(c-\delta)n,
\]
which would contradict the maximality of $F_0$.

We first establish a lower bound on coloured neighbourhoods. Fix $1\le i\le t$ and a vertex $v\in V(G)$.
Since $\chi$ is a proper edge-colouring, all edges incident to $v$ have distinct
colours.
An edge $vw$ is not counted in $N_{C_i}(v)$ only if its colour lies in
$(C(F_0)\cup C_0)\setminus C_i$.
Therefore,
\[
|N_{C_i}(v)|
\ge |N(v)| - |(C(F_0)\cup C_0)\setminus C_i|.
\]
Because $F_0$ is spanning, we have $V(F_0)=V(G)$.
Moreover, $C_0\subset C_i$ and $C_0\cap C(F_0)=\varnothing$, so
\[
|(C(F_0)\cup C_0)\setminus C_i|
=|C(F_0)|-|C_i\setminus C_0|.
\]
Using $\delta(G)\ge cn$ and $|C(F_0)|\le(c-\delta)n$, we obtain
\begin{equation}\label{eq:neighbour-lower}
|N_{C_i}(v)|\ge cn-|C(F_0)|+|C_i\setminus C_0|
\ge \delta n+|C_i\setminus C_0|.
\end{equation}

We next show the growth of the colour sets $C_i\setminus C_0$.
Fix $1\le i\le t$.
By the definition of $C_i$, we have
\[
C_i\setminus C_0
\supseteq
\Big\{f(x):x\in N_{C_{i-1}}(v_1^{(i)})\setminus\{v_1^{(1)},\ldots,v_1^{(t)}\}\Big\}.
\]
By Lemma~\ref{lem:keyobs}, all colours $f(x)$ appearing in the above set are distinct.
Hence
\begin{equation}\label{eq:colour-growth}
|C_i\setminus C_0|
\ge |N_{C_{i-1}}(v_1^{(i)})|-t.
\end{equation}

Applying \eqref{eq:neighbour-lower} with $v=v_1^{(i)}$ and $i-1$ in place of $i$, we get
\[
|N_{C_{i-1}}(v_1^{(i)})|
\ge |C_{i-1}\setminus C_0|+\delta n.
\]
Combining this with \eqref{eq:colour-growth} yields the recurrence
\begin{equation}\label{eq:recurrence}
|C_i\setminus C_0|
\ge |C_{i-1}\setminus C_0|+(\delta n-t)
\qquad\text{for all }1\le i\le t.
\end{equation}

Summing \eqref{eq:recurrence} over $i=1,\ldots,t$, we obtain
\[
|C_t\setminus C_0|
=\sum_{i=1}^t\big(|C_i\setminus C_0|-|C_{i-1}\setminus C_0|\big)
\ge t(\delta n-t).
\]
By assumption, $t<\delta n/2$, and hence $\delta n-t>\delta n/2$.
Together with $t\delta>2c$, this implies
\[
|C_t\setminus C_0|
> \frac{t\delta n}{2}
> (c-\delta)n,
\]
contradicting the assumption that $F_0$ uses fewer than $(c-\delta)n$ colours.

This contradiction completes the proof.
\end{proof}

\end{document}
